% react_barrier: potential energy of two nuclei against centre separation (Vol III Ch 4, fig:react-barrier)
% Units: r in units of the contact distance R1 + R2, U in units of the barrier
% height U_B = q1 q2 / (4 pi eps0 (R1 + R2)).
%   r > 1 :  U = 1/r            (Coulomb repulsion)
%   r < 1 :  U = -1.6           (schematic square well of the strong force)
% A particle of energy E = 0.15 U_B (drawn about seven times below the top; in the
% Sun's core the typical thermal energy is several hundred times below it, so
% the drawing is not to scale) is turned back classically at r = 1/E = 6.67.
% Quantum mechanically it can tunnel from r = 6.67 to r = 1 (dashed).
% Potential energy is plum (book legend); the kinetic-energy level is accent.
\begin{tikzpicture}
\begin{axis}[bookplot, width=10.5cm, height=6.4cm,
    xmin=0, xmax=14.5, ymin=-1.85, ymax=1.3,
    xtick={1}, xticklabels={$R_1+R_2$},
    ytick={0}, yticklabels={$0$},
    xlabel={separation of centres $r$}, ylabel={potential energy $U$},
    ]
  \addplot[zero line,domain=0:14.2,samples=2] {0};
  \fill[cPlum!10] (axis cs:0,-1.6) rectangle (axis cs:1,0);
  \addplot[cPlum,domain=1:14,samples=200] {1/x};
  \addplot[cPlum] coordinates {(0,-1.6) (1,-1.6) (1,1)};
  \draw[leader] (axis cs:0.6,-1.3) -- (axis cs:2.0,-1.3) node[anchor=west,text=cPlum] {strong-force well};
  \node[cPlum,anchor=west] at (axis cs:3.0,0.7) {Coulomb repulsion $U=\dfrac{q_1q_2}{4\pi\varepsilon_0 r}$};
  \draw[leader] (axis cs:1.05,1.0) -- (axis cs:1.6,1.1) node[anchor=west] {barrier height $U_B$};
  % thermal energy and tunnelling
  \addplot[cAcc,thin line,domain=6.667:14,samples=2] {0.15};
  \draw[-smallhead,thin line,cAcc,dash pattern=on 3pt off 2pt] (axis cs:6.667,0.15) -- (axis cs:1.05,0.15);
  \node[cAcc,anchor=south,align=center] at (axis cs:11.5,0.18) {thermal energy};
  \node[cAcc,anchor=north west,align=left] at (axis cs:1.35,-0.04) {tunnelling through\\the barrier};
  \node[dot] at (axis cs:6.667,0.15) {};
  \draw[leader] (axis cs:6.667,0.1) -- (axis cs:8.2,-0.55) node[anchor=north] {classical turning point};
\end{axis}
\end{tikzpicture}
